\section{Теоретическая часть}

\subsection{Оценка среднего значения}

Пусть $x_1, x_2, \dots, x_n$ --- измеренные значения случайной величины. Естественной точечной оценкой математического ожидания является среднее арифметическое независимых измерений:
\[
	\tilde{x}_n = \frac{1}{n} \sum_{i=1}^{n} x_i .
\]

Для случайного процесса $x(t)$ используют дискретные равноотстоящие отсчёты $x(t_i)$, разделённые интервалом $\Delta t$:
\[
	\tilde{x} = \frac{1}{N} \sum_{i=1}^{N} x(t_i), \qquad \Delta t = t_{i+1} - t_i .
\]

При выполнении условий стационарности и эргодичности такая оценка является несмещённой: её статистическое среднее равно математическому ожиданию процесса. При увеличении объёма независимого статистического материала оценка сходится по вероятности к истинному среднему (закон больших чисел).

\subsection{Погрешность оценки и влияние корреляции}

Погрешность оценки характеризуется среднеквадратическим отклонением самой оценки от истинного среднего:
\[
	\sigma_{\tilde{x}} = \sqrt{D[\tilde{x}]} = \sqrt{\big\langle (\tilde{x} - \langle x \rangle)^2 \big\rangle} .
\]

Для $N$ независимых одинаково распределённых отсчётов дисперсия среднего уменьшается в $N$ раз, а СКО --- в $\sqrt{N}$ раз:
\[
	D[\tilde{x}] = \frac{D[x]}{N}, \qquad \sigma_{\tilde{x}} = \frac{\sigma_x}{\sqrt{N}} .
\]

Для коррелированного случайного процесса точность зависит не только от $N$, но и от $\Delta t$ относительно времени корреляции $\tau_\text{корр}$. В методическом пособии дисперсия оценки записана через ковариационную функцию:
\[
	D[\tilde{x}] = \frac{D[x]}{N} + \frac{1}{N^2} \sum_{i=1}^{N} \sum_{j \neq i} B'_x(t_i - t_j) .
\]

Если $\Delta t$ намного меньше $\tau_\text{корр}$, соседние отсчёты сильно зависимы и добавление новых измерений даёт небольшой прирост независимой информации. При $\Delta t$, сравнимом с $\tau_\text{корр}$ или превышающем его, взаимная ковариация соседних отсчётов становится малой, и поведение оценки приближается к случаю независимых измерений.

\subsection{Спектральная плотность мощности}

Корреляционная функция процесса связана с ковариационной функцией и математическим ожиданием соотношением $K_x(\tau) = B_x(\tau) + \langle x \rangle^2$. При ненулевом среднем спектр содержит составляющую на нулевой частоте, связанную с квадратом среднего значения; дисперсия определяется непрерывной частью спектра.

Временное усреднение на интервале $T$ можно представить как линейное преобразование:
\[
	\tilde{x}(t) = \frac{1}{T} \int_{t-T}^{t} x(u) \, du .
\]

Такой усреднитель действует как низкочастотный фильтр. Связь спектров исходного и усреднённого процессов имеет вид:
\[
	S_{\tilde{x}}(2\pi f) = S_x(2\pi f) \, |K(j2\pi f)|^2,
\]
\[
	|K(j2\pi f)|^2 = \left[ \frac{\sin(\pi f T)}{\pi f T} \right]^2 .
\]

Эффективная полоса усреднителя около нулевой частоты имеет порядок $\Delta f_\text{эфф} \approx 1/(2T)$. Поэтому при увеличении $T$ полоса пропускания сужается, а дисперсия оценки уменьшается. Для $T$, значительно большего $\tau_\text{корр}$, число независимых отсчётов внутри интервала усреднения определяет точность оценки.

\subsection{Доверительный интервал}

Поскольку оценка среднего является случайной величиной, её ошибка также случайна. Для заданной доверительной вероятности $\beta$ выбирают доверительный интервал $I_\beta$, который с указанной вероятностью накрывает истинное среднее. При нормальном приближении:
\[
	I_\beta = t_\beta \sigma_{\tilde{x}}, \qquad P\big( |\tilde{x} - \langle x \rangle| \leqslant I_\beta \big) = \beta .
\]

С ростом $\beta$ интервал расширяется; при увеличении $N$ погрешность и доверительный интервал, как правило, уменьшаются.
